\documentclass[a4paper,12pt]{article}
\usepackage[utf8]{inputenc}
\usepackage[indonesian]{babel}
\usepackage[T1]{fontenc}
\usepackage{geometry}
\geometry{margin=2.5cm}
\usepackage{setspace}
\usepackage{hyperref}
\usepackage{graphicx}

\title{Fisiologi Dasar Kardiovaskular}
\author{Renanda Maulana Hamdi}
\date{\today}

\begin{document}

\maketitle
\tableofcontents

\newpage

\section{Pendahuluan}
Sistem kardiovaskular adalah sistem yang terdiri dari jantung, pembuluh darah, dan darah yang memiliki fungsi utama untuk mengangkut bahan-bahan untuk metabolisme jaringan seperti oksigen, dan nutrisi. Kemudian 
membuang produk sisa metabolisme dari jaringan. Selain itu, sistem ini juga berperan dalam regulasi suhu tubuh, pengaturan homeostasis, dan distribusi hormon.
\section{Konsep Dasar Sistem Kardiovaskular}
Sistem kardiovaskular sendiri terdiri dari jantung yang berfungsi sebagai pempompa darah yang berkontraksi dan membuat tekanan untuk menggerakan darah melalui pembuluh darah. Lalu, terdapat pembuluh darah yang berperan sebagai saluran yang mendistribusikan darah ke seluruh jaringan tubuh.
\subsection{Struktur jantung}
\begin{figure}[h]
    \centering
    \includegraphics[width=0.8\textwidth]{.resouces/png/jantung}
    \caption{\href{https://anatomytool.org/content/oli-drawing-heart-and-great-vessels-anterior-english-labels}{"OLI - Drawing Heart and great vessels from anterior - English labels"} by \href{https://oli.cmu.edu/jcourse/lms/students/syllabus.do}{Open Learning Initiative}, license: \href{https://creativecommons.org/licenses/by-nc-sa/4.0/}{CC BY-NC-SA}}
\end{figure}
\begin{itemize}
    \item Jantung terbagi menjadi dua bagian yaitu kanan dan kiri.
    \item Setiap bagian terdiri dari dua ruang yaitu atrium dan ventrikel.
    \item Pada bagian kanan, jantung masuk melalui vena cava superior dan vena cava inferior menuju ke atrium kanan. Lalu, darah keluar dari ventrikel kiri melalui aorta
    \item Pada bagian kiri, jantung masuk melalui vena pulmonalis menuju ke atrium kiri. Lalu, darah keluar dari ventrikel kanan melalui arteri pulmonalis
    \item Diantara ruang-ruang tersebut terdapat katup yang berfungsi sebagai pemisah antar ruang dan agar darah dapat mengalir searah dari atrium ke ventrikel
    \item Katup diantara atrium dan ventrikel disebut sebagai katup atrioventrikular (atrio = atrium, ventrikular = ventrikel)
    \item Katup atrioventrikular pada sisi kanan disebut sebagai katup tricuspid (tri=tiga, cuspis=kelopak, sayap; Karena katup tricuspid memiliki tiga kelopak),  Pada sisi kiri disebut sebagai katup mitral (karena berbentuk seperti topi mitral yang dikenakan paus)
    \item Selain itu, terdapat katup yang memisahkan antara vetrikel dengan arteri yang keluar dari jantung. Katup ini disebut sebagai katup semilunar (semi = separuh,  lunar = bulan; karena berbentuk seperti bulan separuh)
    \item Katup semilunar yang memisahkan ventrikel dengan aorta disebut katup aorta
    \item Katup semilunar yang memisahkan ventrikel dengan arteri/truncus pulmonalis disebut katup pulmonalis
    \item [\textbf{Notes:}] Pembuluh darah yang masuk ke dalam jantung $\rightarrow$ vena, sedangkan pembuluh darah yang keluar dari jantung $\rightarrow$ arteri 
\end{itemize}
\newpage
\subsection{Struktur pembuluh darah}
\begin{figure}[h]
    \centering
    \includegraphics[width=0.8\textwidth]{.resouces/png/arteri et vena.png}
\end{figure}
\section{Fisiologi Jantung}
\begin{itemize}
    \item Siklus jantung: diastole dan sistole
    \item Konduksi listrik jantung
    \item Regulasi denyut jantung
\end{itemize}

\section{Fisiologi Pembuluh Darah dan Sirkulasi}
\begin{itemize}
    \item Anatomi arteri, vena, dan kapiler
    \item Mekanisme kontrol resistensi vaskular
    \item Pertukaran materi di kapiler
\end{itemize}

\section{Regulasi Sistem Kardiovaskular}
\begin{itemize}
    \item Peran sistem saraf otonom
    \item Faktor hormonal dan parakrin
    \item Homeostasis tekanan darah
\end{itemize}

\section{Ringkasan dan Kesimpulan}
Rangkum poin-poin utama, hubungan antar konsep, dan implikasi klinis dasar.

\section{Daftar Pustaka}
\begin{itemize}
    \item Sumber 1
    \item Sumber 2
\end{itemize}

\end{document}